% Part: normal-modal-logic
% Chapter: syntax-and-semantics
% Section: relational-models

\documentclass[../../../include/open-logic-section]{subfiles}

\begin{document}

\olfileid[hi]{nml}{syn}{rel}

\olsection{संबंधात्मक प्रतिरूप}

सामान्य मोडल तर्कशास्त्रों के अर्थविज्ञान की मूल धारणा
\emph{संबंधात्मक प्रतिरूप} है। इसमें द्विपदी ``अभिगम्यता संबंध'' से जुड़े
जगतों का एक समुच्चय और एक ऐसा नियतन होता है जो निर्धारित करता है कि कौन-से
!!{propositional variable} किन जगतों में ``सत्य'' माने जाते हैं।

\begin{defn}
  मूल मोडल भाषा का \emph{प्रतिरूप} एक त्रिक $\mModel{M}
  = \tuple{W, R, V}$ है, जहाँ
  \begin{enumerate}
  \item $W$ ``जगतों'' का अरिक्त समुच्चय है,
  \item $R$, $W$ पर द्विपदी अभिगम्यता संबंध है, और
  \item $V$ ऐसा फलन है जो प्रत्येक !!{propositional variable}~$p$ को संभव जगतों
    का समुच्चय $V(p)$ देता है।
  \end{enumerate}
  जब $Rww'$ मान्य हो, तो हम कहते हैं कि $w'$, $w$ से \emph{अभिगम्य} है।
  जब $w \in V(p)$ हो, तो हम कहते हैं कि $p$, $w$ पर \emph{सत्य} है।
\end{defn}

संबंधात्मक अर्थविज्ञान का बड़ा लाभ यह है कि प्रतिरूपों को
\olref{fig:simple} जैसे सरल आरेखों से निरूपित किया जा सकता है। जगतों को
नोडों से निरूपित किया जाता है, और जगत~$w'$, $w$ से ठीक तभी अभिगम्य है जब
$w$ से $w'$ की ओर कोई तीर हो। इसके अतिरिक्त, $w \in V(p)$ होने पर हम किसी
नोड (जगत) पर~$\mTrue{p}$ और अन्यथा~$\mFalse{p}$ अंकित करते हैं।
\olref{fig:simple} उस प्रतिरूप को निरूपित करता है जिसमें
$W = \{w_1, w_2, w_3\}$, $R = \{\tuple{w_1, w_2},\tuple{w_1,
  w_3}\}$, $V(p) = \{w_1, w_2\}$ और $V(q) = \{w_2\}$।

\begin{figure}
  \begin{center}
    \begin{tikzpicture}[modal]
      \node[world] (w1) [label={[align=right]right:\mTrue{p}\\ \mFalse{q}}]{$w_1$}; 
      \node[world] (w2) [label={[align=right]right:\mTrue{p}\\ \mTrue{q}},
        above right=of w1]{$w_2$}; 
      \node[world] (w3) [label={[align=right]right:\mFalse{p}\\ \mFalse{q}},
        below right=of w1] {$w_3$};
      \draw[->] (w1) to (w2);
      \draw[->] (w1) to (w3);
    \end{tikzpicture}
  \end{center}
  \caption{एक सरल प्रतिरूप।}
  \ollabel{fig:simple}
\end{figure}

\end{document}
